\chapter{Хопфово-скрученный дефект: совместный тест энергии и индекса}

Предыдущий гейт канонически назначил хопфовы линии $L/L^*$ направленным
стрелкам $E/E^*$. Тем самым единичный первый класс Черна больше не является
внешним знаком. Остаётся более сильный вопрос: может ли одна и та же
геометрия одновременно сделать точечный дефект конечным по энергии и дать
фермионный индекс один.

Настоящий гейт намеренно различает три утверждения:
\begin{enumerate}
 \item что даёт одна абелева хопфова связность;
 \item существует ли гладкое математическое завершение с тем же зарядом;
 \item выведено ли это завершение из текущего родителя $M_{35}$.
\end{enumerate}
Смешение этих уровней превратило бы стандартный монопольный пример в
ложное доказательство нашей конкретной теории.

\section{Почему одной хопфовой линии недостаточно}

На сфере вокруг дефекта хопфова линия имеет
\begin{equation}
 \frac1{2\pi}\int_{S^2}F_L=1.
 \label{eq:v5-hopf-energy-index-c1}
\end{equation}
Если продолжить эту абелеву кривизну внутрь как поле точечного магнитного
заряда, то при нормировке потока $2\pi$ получаем
\begin{equation}
 B_r=\frac1{2r^2}.
\end{equation}
Её максвелловская энергия вблизи начала содержит
\begin{equation}
 \frac12\int_0^\varepsilon 4\pi r^2 B_r^2\,dr
 =\frac\pi2\int_0^\varepsilon\frac{dr}{r^2},
 \label{eq:v5-hopf-dirac-core-divergence}
\end{equation}
то есть расходится на ядре.

С другой стороны, абелева фаза действует вертикально на хопфов спинор
$z\mapsto e^{i\alpha}z$ и оставляет
\begin{equation}
 n=z^\dagger\boldsymbol\sigma z,
 \qquad
 P=nn^T
\end{equation}
неизменными. Поэтому такая $U(1)$-связность не может компенсировать
угловой градиент самого проекторного ежа. Для него по-прежнему остаётся
оболочечная энергия $8\pi\,dr$ и линейная расходимость на бесконечности.

\begin{nogocriterion}[Абелева линия не является полным ядром частицы]
Хопфова линия с $c_1=1$ является правильным носителем ориентации и
индекса, но её чистая $U(1)$-кривизна не даёт конечную энергию: она не
экранирует проекторный градиент и сингулярна при точечном продолжении.
\end{nogocriterion}

Это важное отрицательное уточнение. Слово ``одна'' в условии предыдущего
гейта нельзя понимать как ``только абелева''. Нужна одна гладкая
неабелева связность, чьим асимптотическим собственным подрасслоением
является уже выведенная линия $L$.

\section{Минимальное гладкое завершение}

Канонический математический образец даёт заряд-один монополь
Богомольного--Прасада--Соммерфилда. Пусть $A_i^a$ --- пространственная
$SO(3)$-связность, а $\Phi^a$ --- поле в присоединённом представлении.
Для $\rho=gvr$ возьмём сферически-симметричный анзац
\begin{align}
 A_i^a&=\epsilon_{aij}\frac{x^j}{gr^2}\bigl(1-w(\rho)\bigr),\\
 \Phi^a&=v\frac{x^a}{r}h(\rho),
 \label{eq:v5-hopf-bps-ansatz}
\end{align}
с профилями
\begin{equation}
 w(\rho)=\frac{\rho}{\sinh\rho},
 \qquad
 h(\rho)=\coth\rho-\frac1\rho.
 \label{eq:v5-hopf-bps-profiles}
\end{equation}
Они удовлетворяют
\begin{equation}
 w'=-wh,
 \qquad
 h'=\frac{1-w^2}{\rho^2},
 \label{eq:v5-hopf-bps-odes}
\end{equation}
и граничным условиям
\begin{equation}
 w(0)=1,quad h(0)=0,qquad
 w(\infty)=0,quad h(\infty)=1.
 \label{eq:v5-hopf-bps-boundaries}
\end{equation}

В начале координат связность и поле порядка гладко распускают точечную
сингулярность. На бесконечности направление $\widehat\Phi$ совпадает с
ежом $n=x/r$, а неабелева связность компенсирует его угловое изменение.
Тем самым один объект устраняет обе расходимости чистой абелевой модели.

Для функционала Янга--Миллса--Хиггса в пределе Богомольного
\begin{equation}
 E=\frac12\int_{\mathbb R^3}
 \left(|B|^2+|D_A\Phi|^2\right)d^3x
\end{equation}
выполнено разложение
\begin{equation}
 E=\frac12\int|B-D_A\Phi|^2d^3x
 +\int_{S^2_\infty}\!\langle\Phi,B\rangle.
 \label{eq:v5-hopf-bogomolny-square}
\end{equation}
Решение $B=D_A\Phi$ имеет заряд один и конечную энергию
\begin{equation}
 E_{\rm BPS}=\frac{4\pi v}{g}.
 \label{eq:v5-hopf-bps-energy}
\end{equation}
Формулы профиля были впервые получены Прасадом и Соммерфилдом;
аналитический индекс на открытом пространстве задаётся теоремой Каллиаса.

\section{Тот же заряд как фермионный индекс}

Рассмотрим скрученный оператор
\begin{equation}
 \mathcal D_{A,\Phi}=\sigma^i(\partial_i+A_i)+\Phi.
 \label{eq:v5-hopf-callias-operator}
\end{equation}
На сфере бесконечности $\Phi$ невырождено, и его положительное собственное
подрасслоение является именно хопфовой линией $L$:
\begin{equation}
 P_+(n)=\frac12(I+n\cdot\boldsymbol\sigma),
 \qquad c_1(L)=1.
\end{equation}
Для фундаментального комплексного фермиона индекс Каллиаса равен этому
магнитному заряду,
\begin{equation}
 \operatorname{ind}\mathcal D_{A,\Phi}=1.
 \label{eq:v5-hopf-callias-index-one}
\end{equation}

Следовательно, энергия и индекс здесь не чинятся двумя независимыми
добавками. Их связывает одна пара $(A,\Phi)$ и один граничный класс:
\begin{equation}
 \boxed{
 c_1(L)=Q_{\rm mag}=1
 \quad\Longrightarrow\quad
 E_{\rm BPS}=4\pi v/g,
 \quad
 \operatorname{ind}\mathcal D_{A,\Phi}=1.}
 \label{eq:v5-hopf-joint-charge}
\end{equation}

Это не доказывает, что физический нейтринный сектор проекта уже реализует
фундаментальный модуль данного $SU(2)$. Утверждение пока имеет статус
математического совместного образца: нужные два эффекта принципиально
совместимы и управляются тем самым классом, который был получен из
ориентационного функтора Мориты.

\section{Суперсвязностное чтение}

На градуированном пакете
\begin{equation}
 \mathcal E_H=(E\otimes L)\oplus(E^*\otimes L^*)
\end{equation}
пару калибровочного и нечётного полей естественно записать как
суперсвязность
\begin{equation}
 \mathbb A=\nabla_A+\Phi.
\end{equation}
Её кривизна содержит совместно
\begin{equation}
 \mathbb F=F_A+D_A\Phi+\Phi^2.
 \label{eq:v5-hopf-supercurvature}
\end{equation}
Именно эта упаковка соответствует языку переходов проекта: диагональная
кривизна, ковариантная производная перехода и квадрат перехода являются
частями одного градуированного объекта.

KO6 обменивает две ориентированные ветви
\begin{equation}
 E\otimes L\longleftrightarrow E^*\otimes L^*.
\end{equation}
Поэтому полный вещественный пакет содержит сопряжённые индексы $+1$ и
$-1$. Индекс один относится к выбранной ориентированной физической
ветви; выдавать сумму полного KO6-пакета за ненулевой комплексный индекс
нельзя.

\section{Что уже есть в проекте}

Гладкое завершение не появляется из пустоты. Ранний гейт
\path{family_connection_defect_gap_bridge.tex} уже построил семейную
$SO(3)$-связность на трубчатой окрестности дефекта и показал, что она
может связывать объёмное семейное расщепление с локализованным ядром.
Текущая ветвь добавляет к этому два существенных результата:
\begin{enumerate}
 \item ориентация связи больше не требует нового фиксированного
 $SU(2)_F$-дублета;
 \item её граничная собственная линия канонически назначена стрелке
 $E$ и имеет $c_1=1$.
\end{enumerate}

Однако прежняя семейная связность была локальным условным мостом. Она не
была выведена как глобальное пространственное поле Янга--Миллса--Хиггса и
не сопровождалась единственным радиальным профилем
\eqref{eq:v5-hopf-bps-profiles}.

\section{Невыведенные данные родителя}

Чтобы формула \eqref{eq:v5-hopf-joint-charge} стала результатом версии V,
из одного родительского принципа должны следовать:
\begin{enumerate}
 \item глобальная пространственная $SO(3)$-связность $A$;
 \item амплитуда $|\Phi|$ и её регулярное обращение в нуль на ядре;
 \item относительный коэффициент между $|F_A|^2$ и $|D_A\Phi|^2$;
 \item связь фундаментального оператора Каллиаса с физическим пакетом
 $H_{15}$;
 \item масштабы $g$ и $v$ либо правило, отделяющее топологический вывод
 от их последующего измерения.
\end{enumerate}

Нормированный конечный след $M_{35}$ сам по себе фиксирует внутренние
матричные следы, но ещё не выводит пространственную меру, жёсткость
калибровочного поля или размерный масштаб. Поэтому подстановка стандартного
BPS-действия сейчас была бы внешним завершением, а не выводом.

\begin{theorem}[Совместимость энергии и индекса]
Хопфов класс $c_1=1$, канонически назначенный ориентированному переходу,
допускает гладкое неабелево BPS-завершение. В нём одна пара $(A,\Phi)$
имеет конечную энергию и даёт индекс Каллиаса один. Следовательно,
совместный энергетико-индексный механизм математически существует.
Текущий родитель $M_{35}$ ещё не выводит его пространственное действие и
нормировку.
\end{theorem}

\section{Вердикт и следующий гейт}

Мы не пришли ни к новому тупику, ни к готовой частице. Получен более узкий
и полезный результат:
\begin{equation}
 \boxed{
 \begin{gathered}
 \text{чистая хопфова }U(1)\text{-линия: индекс да, энергия нет},\\
 \text{гладкое }SO(3)\text{-завершение: энергия и индекс да},\\
 \text{вывод этого завершения из }M_{35}:\text{ пока нет}.
 \end{gathered}}
\end{equation}

Следующий допустимый вопрос теперь однозначен: порождает ли градуированная
суперсвязность Мориты глобальную пространственную $SO(3)$-кривизну и
богомольный относительный вес из одного следа, или для этого неизбежно
нужно новое действие. Этому посвящается гейт
\path{version5_spatial_so3_superconnection_parent_trace_gate.tex}.

\begin{protocol}
\begin{enumerate}
 \item чистая хопфова линия имеет $c_1=1$: \textbf{да};
 \item её абелева энергия конечна: \textbf{нет};
 \item она компенсирует градиент $P$: \textbf{нет};
 \item гладкое неабелево завершение существует: \textbf{да};
 \item BPS-профили проходят уравнения первого порядка: \textbf{да};
 \item энергия завершения конечна: \textbf{$4\pi v/g$};
 \item заряд завершения: \textbf{один};
 \item индекс фундаментальной ориентированной ветви: \textbf{один};
 \item энергия и индекс управляются одним зарядом: \textbf{да};
 \item KO6-сопряжённая ветвь сохранена: \textbf{да};
 \item глобальное пространственное действие выведено из $M_{35}$:
 \textbf{нет};
 \item масштабы и жёсткость выведены: \textbf{нет};
 \item математический совместный тест: \textbf{пройден};
 \item физическое замыкание версии V: \textbf{не достигнуто}.
\end{enumerate}
\end{protocol}

Машинный протокол сохранён в
\path{s2t_v5_hopf_twisted_defect_superconnection_energy_index_gate_results.json}.